\chapter{Duality, Complementary Slackness, and the Price Theorem}

Every linear program comes with a hidden companion: its \emph{dual}. Where the primal asks how to use scarce resources to
maximise profit, the dual asks what those resources are actually worth. The
two problems are locked together by a family of results (weak duality,
strong duality, and complementary slackness) that turn the dual solution
into a certificate of optimality for the primal and, at the same time, into
a price list: the optimal dual variables are the \emph{shadow prices} of the constraints.

This chapter drills the mechanics: constructing the dual of a general LP
(including $\ge$ and $=$ constraints and free variables), bounding and
certifying optima with weak duality, recovering one optimal solution from
the other via complementary slackness, and reading the dual solution
straight out of the simplex end tableau. The price
theorem then gives the economic payoff: it tells you,
constraint by constraint, how much an extra unit of each resource is worth.

\section{Core Rules}

\begin{rulebox}{The standard primal--dual pair}
For a primal LP in standard maximum form the dual is built as follows:
\[
  \textbf{Primal:}\;\;
  \max\; Z=c^{\top}x,\;\; Ax\le b,\;\; x\ge 0
  \qquad
  \textbf{Dual:}\;\;
  \min\; W=b^{\top}y,\;\; A^{\top}y\ge c,\;\; y\ge 0 .
\]
Construction recipe:
\begin{itemize}[leftmargin=1.6em,itemsep=0.3ex]
  \item one dual variable $y_{i}$ per primal \emph{constraint}, one dual
        constraint per primal \emph{variable};
  \item the primal right-hand sides $b_{i}$ become the dual objective
        coefficients; the primal objective coefficients $c_{j}$ become the
        dual right-hand sides;
  \item the coefficient matrix is transposed: the \emph{column} of $x_{j}$
        in the primal is the \emph{row} of the $j$-th dual constraint;
  \item $\max\leftrightarrow\min$ and $\le\leftrightarrow\ge$.
\end{itemize}
Applying the recipe twice returns the original problem: \emph{the dual of
the dual is the primal}.
\end{rulebox}

\begin{rulebox}{Sign rules for the general dual}
When the primal is a maximum problem with mixed constraint types:
\begin{center}
\begin{tabular}{l|l}
\textbf{Primal (max)} & \textbf{Dual (min)}\\
\hline
$i$-th constraint $\le b_{i}$ & $y_{i}\ge 0$\\
$i$-th constraint $\ge b_{i}$ & $y_{i}\le 0$\\
$i$-th constraint $= b_{i}$   & $y_{i}$ free\\
\hline
variable $x_{j}\ge 0$ & $j$-th constraint $\ge c_{j}$\\
variable $x_{j}\le 0$ & $j$-th constraint $\le c_{j}$\\
variable $x_{j}$ free & $j$-th constraint $= c_{j}$\\
\end{tabular}
\end{center}
The table is symmetric: for a primal in \emph{minimum} form, locate its
features in the right-hand column and read off the dual feature on the
left. (Example: a $\ge$ constraint of a min problem gets a dual variable
$u_{i}\ge 0$; a variable $x_{j}\ge 0$ of a min problem gets a dual
constraint $\le c_{j}$.) Memory aid: in a max problem, $\le$ constraints
and $\ge 0$ variables are the ``natural'' ingredients and produce
sign-restricted partners; each ``unnatural'' ingredient flips the sign,
and equalities pair with free variables.
\end{rulebox}

\begin{rulebox}{Weak duality}
For \emph{every} primal-feasible $x$ and \emph{every} dual-feasible $y$ of
a max/min pair:
\[
  Z(x) \;=\; c^{\top}x \;\le\; b^{\top}y \;=\; W(y) .
\]
Consequences:
\begin{itemize}[leftmargin=1.6em,itemsep=0.3ex]
  \item any dual-feasible $y$ gives an \emph{upper bound} on $Z^{*}$; any
        primal-feasible $x$ gives a \emph{lower bound} on $W^{*}$;
  \item \textbf{certificate:} if a feasible pair satisfies
        $Z(x)=W(y)$, then $x$ is primal-optimal and $y$ is dual-optimal;
  \item if the primal is unbounded, the dual is infeasible (and vice
        versa): otherwise a dual-feasible point would cap $Z$.
\end{itemize}
\end{rulebox}

\begin{rulebox}{Strong duality}
If the primal has an optimal solution $x^{*}$, then the dual has an
optimal solution $y^{*}$ and the optimal values coincide:
\[
  Z^{*} \;=\; c^{\top}x^{*} \;=\; b^{\top}y^{*} \;=\; W^{*} .
\]
The bound from weak duality is therefore \emph{tight}: solving either
problem solves both.
\end{rulebox}

\begin{rulebox}{Complementary slackness}
Feasible $x^{*}$ and $y^{*}$ are optimal for their respective problems if
and only if
\[
  y_{i}^{*}\,\bigl(b_{i}-a_{i}^{\top}x^{*}\bigr)=0 \;\;\forall i,
  \qquad
  x_{j}^{*}\,\bigl((A^{\top}y^{*})_{j}-c_{j}\bigr)=0 \;\;\forall j .
\]
In words: a slack primal constraint forces its dual variable to zero, and
a positive primal variable forces its dual constraint to bind, and
symmetrically in the other direction. Recipe for recovering $y^{*}$ from a
known primal optimum $x^{*}$:
\begin{itemize}[leftmargin=1.6em,itemsep=0.3ex]
  \item constraint $i$ has slack at $x^{*}$ $\;\Rightarrow\;$ set
        $y_{i}=0$;
  \item $x_{j}^{*}>0$ $\;\Rightarrow\;$ write the $j$-th dual constraint
        as an \emph{equation};
  \item solve the resulting linear system, then check dual feasibility and
        $W=Z^{*}$.
\end{itemize}
The same recipe run backwards recovers $x^{*}$ from a known $y^{*}$.
\end{rulebox}

\begin{rulebox}{The price theorem (shadow prices)}
The optimal dual variable of constraint $i$ measures the marginal value of
its right-hand side:
\[
  y_{i}^{*} \;=\; \pd{Z^{*}}{b_{i}} ,
\]
i.e.\ enlarging the resource stock $b_{i}$ by one (small) unit raises the
optimal value by $y_{i}^{*}$: valid as long as the optimal basis does not
change. Hence $y_{i}^{*}$ is the \emph{shadow price} of resource $i$:
\begin{itemize}[leftmargin=1.6em,itemsep=0.3ex]
  \item \textbf{decision rule:} buy an extra unit of resource $i$ if and
        only if its market price is below $y_{i}^{*}$;
  \item a constraint with slack has shadow price $0$: unused capacity is
        worthless at the margin;
  \item only binding constraints can carry a positive shadow price.
\end{itemize}
\end{rulebox}

\begin{rulebox}{The dual as a pricing problem}
Read the dual of a production LP as the problem of an outside bidder who
wants to rent the firm's entire resource stock:
\begin{itemize}[leftmargin=1.6em,itemsep=0.3ex]
  \item $y_{i}$ = offered rental price per unit of resource $i$;
  \item dual constraint $j$: the resources consumed by one unit of product
        $j$, valued at prices $y$, must be worth at least its profit
        $c_{j}$: otherwise the firm would rather produce than rent out;
  \item dual objective: minimise the total rent
        $b^{\top}y$ for the stock;
  \item strong duality: the cheapest acceptable offer equals the maximal
        profit $Z^{*}$ the firm could earn itself.
\end{itemize}
\end{rulebox}

\begin{rulebox}{Dual values in the end tableau}
For a max LP in $\le$ form solved with slack variables $s_{1},\dots,s_{m}$:
\begin{itemize}[leftmargin=1.6em,itemsep=0.3ex]
  \item the bottom-row entry under slack column $s_{i}$ in the end tableau
        is the optimal dual variable $y_{i}^{*}$;
  \item the bottom-right entry is $Z^{*}=W^{*}$;
  \item a slack variable that is basic with positive value marks a
        non-binding constraint, so its bottom-row entry (its shadow price)
        is $0$;
  \item the primal solution sits in the RHS column (basic variables);
        non-basic variables are $0$.
\end{itemize}
\end{rulebox}

\begin{rulebox}{Dual slack variables}
Writing the $j$-th dual constraint of the standard pair as an equation
introduces a \emph{dual slack} (surplus)
\[
  t_{j} \;=\; (A^{\top}y)_{j}-c_{j}\;\ge 0,
\]
that is, the opportunity cost of the resources one unit of product~$j$
uses, minus its margin.
At the optimum $t_{j}$ is the $z$-row entry under $x_{j}$ in the end
tableau, i.e.\ the absolute value of the reduced cost of Chapter~8, and
complementary slackness reads $t_{j}\,x_{j}^{*}=0$. Hence a product that is
made has $t_{j}=0$ (it exactly breaks even at shadow prices), and $t_{j}>0$
means the product is not made: its margin would have to rise by more than
$t_{j}$ before producing it pays.
\end{rulebox}

\begin{rulebox}{Complementary slackness with sign-restricted or free duals}
The recipe is the same for every primal: a row with slack gets a zero dual
variable, a positive primal variable turns its dual constraint into an
equation. After solving, check each dual variable against the sign rules:
\begin{itemize}[leftmargin=1.6em,itemsep=0.3ex]
  \item \textbf{max primal:} $\le$ row $y_{i}\ge 0$; $\ge$ row
        $y_{i}\le 0$; $=$ row $y_{i}$ free. A binding minimum requirement
        ($\ge$) has $y_{i}\le 0$: raising the requirement by one unit
        \emph{lowers} $Z^{*}$ by $|y_{i}|$;
  \item \textbf{min primal:} $\ge$ row $u_{i}\ge 0$ (a stricter requirement
        raises the cost); $\le$ row $u_{i}\le 0$ (more capacity lowers the
        cost); $=$ row $u_{i}$ free. In each case $u_{i}=\partial C^{*}/\partial b_{i}$.
\end{itemize}
Converting a row by $-1$ (for example $a^{\top}x\ge b$ into
$-a^{\top}x\le -b$) flips the sign of its dual variable and nothing else:
the converted row of a max problem has the price $-y_{i}\ge 0$. A dual
value with the wrong sign means an error, or a wrong guess of the binding
set.
\end{rulebox}

\begin{rulebox}{Solving the dual instead of the primal}
Pick the smaller problem. A primal with $n\ge 3$ variables but only two
constraints has a dual with two variables, which can be solved
graphically:
\begin{enumerate}[leftmargin=1.8em,itemsep=0.3ex]
  \item draw the $n$ dual constraints $(A^{\top}y)_{j}\ge c_{j}$ in the
        $(y_{1},y_{2})$-plane, find the optimal corner $y^{*}$ of
        $\min b^{\top}y$ and the value $W^{*}$;
  \item a dual constraint that is slack at $y^{*}$ forces $x_{j}^{*}=0$
        (that product is not made);
  \item each $y_{i}^{*}>0$ forces primal row $i$ to bind;
  \item solve the resulting (usually $2\times2$) system for the remaining
        $x_{j}$, check $x\ge 0$ and $Z^{*}=W^{*}$.
\end{enumerate}
Read-out: products made are those whose dual constraint binds; a resource
with $y_{i}^{*}=0$ may be left over (compute its slack); a product not made
needs a margin of at least $(A^{\top}y^{*})_{j}$, i.e.\ $t_{j}$ more than
it has. Conversely, a min primal with $\ge$ rows and cost coefficients
$c\ge 0$ has a dual $\max b^{\top}u$, $A^{\top}u\le c$, $u\ge 0$ that is a
standard maximum problem: solve the dual by the plain
simplex method, and read the primal optimum from the end $z$-row under the
dual's slack columns, while the entries under the dual's own variables are
the primal surpluses.
\end{rulebox}

\begin{rulebox}{LP duality and the Kuhn--Tucker conditions}
For $\max c^{\top}x$ s.t.\ $Ax\le b$, $x\ge 0$ form the Lagrangian with one
multiplier per resource row:
\[
  \mathcal{L}(x,y)=c^{\top}x+y^{\top}(b-Ax)
  \;=\;b^{\top}y+x^{\top}(c-A^{\top}y).
\]
The Kuhn--Tucker conditions (with $x\ge 0$) are
\[
  \begin{aligned}
    &\pd{\mathcal{L}}{x_{j}}=c_{j}-(A^{\top}y)_{j}\le 0,\quad x_{j}\ge 0,\quad
     x_{j}\pd{\mathcal{L}}{x_{j}}=0,\\
    &\pd{\mathcal{L}}{y_{i}}=b_{i}-(Ax)_{i}\ge 0,\quad y_{i}\ge 0,\quad
     y_{i}\pd{\mathcal{L}}{y_{i}}=0 .
  \end{aligned}
\]
Line by line: the first group is \emph{dual feasibility} $A^{\top}y\ge c$
plus the complementary slackness $x_{j}t_{j}=0$; the second group is
\emph{primal feasibility} plus $y_{i}s_{i}=0$. So the Kuhn--Tucker
multipliers are exactly the dual variables, the multiplier of $x_{j}\ge 0$
(in the version with explicit $\mu_{j}$) is the dual slack $t_{j}$, and the
multiplier of a resource is its shadow price. Because the objective and the
constraints are linear (hence concave and convex), the Kuhn--Tucker
conditions are not only necessary but also sufficient: a pair $(x,y)$
satisfying them is an optimal primal--dual pair.
\end{rulebox}

\section{Worked Example}

\begin{worked}{Dualise the running example, solve the dual, and verify strong duality and complementary slackness}
\wstep{Step 1. Write the primal in the standard maximum form.}
\[
  \begin{array}{lrcrcl}
    \max & 40x_{1} & + & 30x_{2}\\
    \text{s.t.} & 2x_{1} & + & x_{2} & \le & 100\\
                & x_{1} & + & x_{2} & \le & 80\\
                & \multicolumn{5}{l}{x_{1},x_{2}\ge0}
  \end{array}
\]

\wstep{Step 2. Transpose everything.}
The dual has \textbf{one variable per primal constraint} and \textbf{one
constraint per primal variable}. Right-hand sides and objective
coefficients swap places, and the matrix is transposed:
\[
  \begin{array}{lrcrcl}
    \min & 100y_{1} & + & 80y_{2}\\
    \text{s.t.} & 2y_{1} & + & y_{2} & \ge & 40 \quad(\text{from }x_{1})\\
                & y_{1} & + & y_{2} & \ge & 30 \quad(\text{from }x_{2})\\
                & \multicolumn{5}{l}{y_{1},y_{2}\ge0}
  \end{array}
\]
\wtrap{Everything flips: max becomes min, $\le$ becomes $\ge$, the objective
coefficients $(40,30)$ become right-hand sides, and the right-hand sides
$(100,80)$ become objective coefficients. The one thing that does not flip
is the sign restriction $y\ge0$, which comes from the primal's $\le$
constraints. Build the dual by transposing the whole tableau of numbers, not
by translating the words.}

\wstep{Step 3. Read the dual's meaning before solving it.}
$y_{i}$ is a price per unit of resource~$i$. The first dual constraint says:
the resources consumed by one unit of product~1 must be valued at no less
than the $40$ euros that unit earns. The dual minimises the total imputed
value of the endowment subject to no product being undervalued --- it is the
problem of an accountant pricing the firm's resources, and the primal is the
problem of the production manager using them.

\wstep{Step 4. Solve the dual.}
Both primal variables are positive at the optimum, so (by Step~7 below) both
dual constraints hold with equality:
\[
  2y_{1}+y_{2}=40,\qquad y_{1}+y_{2}=30
  \;\Longrightarrow\;
  y_{1}=10,\ y_{2}=20 .
\]
Both are $\ge0$ \checkmark, and the dual objective is
\[
  100(10)+80(20)=1000+1600=2600 .
\]

\wstep{Step 5. Strong duality, verified.}
\[
  \boxed{z^{*}_{\text{primal}}=2600=w^{*}_{\text{dual}} .}
\]
The two problems have the \emph{same} optimal value. And these are exactly
the numbers standing in the $z$-row of the end tableau in Chapter~5 --- 
solving the primal by simplex solves the dual for free.

\wstep{Step 6. Weak duality, and what it is for.}
For \emph{any} feasible $\mathbf x$ and \emph{any} feasible $\mathbf y$,
\[
  \mathbf c^{\!\top}\mathbf x\le\mathbf b^{\!\top}\mathbf y .
\]
Every feasible dual solution therefore certifies an \textbf{upper bound} on
the primal, and every feasible primal solution a \textbf{lower bound} on the
dual. Take $\mathbf y=(15,15)$: it satisfies $2(15)+15=45\ge40$ and
$15+15=30\ge30$, so it is dual-feasible, and it bounds the primal by
\[
  100(15)+80(15)=2700 .
\]
Take the primal point $(0,80)$: feasible, worth $2400$. So without solving
anything we know the optimum lies in $[2400,2700]$ --- and the true value
$2600$ does. \checkmark
\wtrap{The bound runs one way only. If a feasible dual value ever came out
\emph{below} a feasible primal value, one of the two points is not actually
feasible or the arithmetic is wrong --- weak duality makes that
impossible. This is the cheapest error check in the volume: evaluate both
objectives and confirm the primal value is the smaller.}

\wstep{Step 7. Complementary slackness, term by term.}
At the optimum, for each pair:
\[
  y_{i}\cdot(\text{slack in primal constraint }i)=0,
  \qquad
  x_{j}\cdot(\text{surplus in dual constraint }j)=0 .
\]
\begin{center}\footnotesize
\begin{tabular}{lll}
  \hline
  pair & values & product\\
  \hline
  $y_{1}=10$, machine slack $=0$ & $10\times0$ & $0$ \checkmark\\
  $y_{2}=20$, labour slack $=0$ & $20\times0$ & $0$ \checkmark\\
  $x_{1}=20$, dual surplus $2(10)+20-40=0$ & $20\times0$ & $0$ \checkmark\\
  $x_{2}=60$, dual surplus $10+20-30=0$ & $60\times0$ & $0$ \checkmark\\
  \hline
\end{tabular}
\end{center}
In words: \emph{a resource with spare capacity is free, and a product that
is made must break even at shadow prices}.

\wstep{Step 8. Use complementary slackness as a solution method.}
Given the primal solution $(20,60)$, both variables are positive, so both
dual constraints must be tight --- which is exactly how Step~4 solved a
two-variable dual with no simplex at all. This works in reverse too: given a
claimed optimal $\mathbf y$, the tight dual constraints identify which
$x_{j}$ may be non-zero, and the binding primal constraints then pin them
down.
\wtrap{Complementary slackness certifies a \emph{pair}. A primal solution
alone, or a dual solution alone, proves nothing; the two together, feasible
and complementary, prove both are optimal. That is the cheapest correctness
proof available in this volume --- and the only one that does not require
re-running the algorithm.}
\end{worked}

\section{Practice Problems}

\begin{problems}
  \item Write the dual of
        \[
          \max\; Z=3x_{1}+2x_{2}
          \quad\text{s.t.}\quad
          x_{1}+x_{2}\le 4,\qquad 2x_{1}+x_{2}\le 6,\qquad
          x_{1},x_{2}\ge 0 .
        \]
  \item Write the dual of
        \[
          \max\; Z=4x_{1}+3x_{2}
          \quad\text{s.t.}\quad
          2x_{1}+x_{2}\le 10,\qquad x_{1}+x_{2}\le 6,\qquad
          x_{1}+2x_{2}\le 10,\qquad x_{1},x_{2}\ge 0 .
        \]
  \item Write the dual of the minimum problem
        \[
          \min\; Z=6x_{1}+4x_{2}
          \quad\text{s.t.}\quad
          3x_{1}+x_{2}\ge 3,\qquad x_{1}+x_{2}\ge 2,\qquad
          x_{1}+2x_{2}\ge 2,\qquad x_{1},x_{2}\ge 0 .
        \]
  \item Using the general sign rules, write the dual of
        \[
          \max\; Z=2x_{1}+5x_{2}
          \quad\text{s.t.}\quad
          x_{1}+3x_{2}\le 12,\qquad x_{1}-x_{2}\ge 2,\qquad
          x_{1}+x_{2}= 8,\qquad x_{1},x_{2}\ge 0 .
        \]
  \item Write the dual of
        \[
          \max\; Z=x_{1}+4x_{2}
          \quad\text{s.t.}\quad
          2x_{1}+x_{2}\le 8,\qquad x_{1}+2x_{2}\le 10,\qquad
          x_{1}\ge 0,\; x_{2}\ \text{free}.
        \]
  \item For the LP
        $\max\; Z=5x_{1}+4x_{2}$ s.t.\ $6x_{1}+4x_{2}\le 24$,
        $x_{1}+2x_{2}\le 6$, $x_{1},x_{2}\ge 0$, form the dual, then form
        the dual of the dual, and confirm that the primal reappears.
  \item Consider the LP of Problem~1. Show that $\hat{x}=(1,2)$ is primal
        feasible and $\hat{y}=(1,1)$ is dual feasible, compute
        $Z(\hat{x})$ and $W(\hat{y})$, verify the weak-duality inequality,
        and state the bounds this pair puts on the optimal value $Z^{*}$.
  \item Still for the LP of Problem~1: show that $x=(2,2)$ and $y=(1,1)$
        are feasible with $Z(x)=W(y)$, and conclude (using weak duality
        only) that both are optimal.
  \item Consider
        $\max\; Z=x_{1}+x_{2}$ s.t.\ $x_{1}-x_{2}\le 1$,
        $x_{1},x_{2}\ge 0$.
        (a)~Show that the LP is unbounded.
        (b)~Write its dual and show that the dual is infeasible.
        (c)~Explain why weak duality makes (b) a necessary consequence of
        (a).
  \item The LP of Problem~2 has the optimal solution $x^{*}=(4,2)$. Use
        complementary slackness to compute the optimal dual solution
        $y^{*}$, and verify strong duality.
  \item Consider
        $\max\; Z=2x_{1}+x_{2}$ s.t.\ $x_{1}+x_{2}\le 4$,
        $x_{1}+3x_{2}\le 6$, $x_{1},x_{2}\ge 0$.
        Its dual has the optimal solution $y^{*}=(2,0)$. Use complementary
        slackness to recover the optimal primal solution $x^{*}$ and
        verify $Z^{*}=W^{*}$.
  \item For the LP of Problem~1, check whether the feasible pair
        $x=(3,0)$, $y=(1,1)$ satisfies all complementary-slackness
        conditions. What do you conclude?
  \item The minimum problem
        $\min\; Z=2x_{1}+3x_{2}$ s.t.\ $x_{1}+x_{2}\ge 5$,
        $x_{1}+2x_{2}\ge 6$, $x_{1},x_{2}\ge 0$
        has the optimal solution $x^{*}=(4,1)$. Write the dual and use
        complementary slackness to find its optimal solution $u^{*}$;
        verify the objective values agree.
  \item Solving
        $\max\; Z=3x_{1}+2x_{2}$ s.t.\
        $x_{1}+x_{2}\le 4$ (resource A, slack $s_{1}$),
        $2x_{1}+x_{2}\le 6$ (resource B, slack $s_{2}$),
        $x_{1},x_{2}\ge 0$
        by the simplex method produces the end tableau
        \begin{center}
        \begin{tabular}{c|cccc|c}
          Basis & $x_{1}$ & $x_{2}$ & $s_{1}$ & $s_{2}$ & RHS\\
          \hline
          $x_{1}$ & $1$ & $0$ & $-1$ & $1$  & $2$\\
          $x_{2}$ & $0$ & $1$ & $2$  & $-1$ & $2$\\
          \hline
          $Z$     & $0$ & $0$ & $1$  & $1$  & $10$\\
        \end{tabular}
        \end{center}
        Read off (a)~the optimal solution and $Z^{*}$, (b)~which
        constraints bind, (c)~the shadow prices, and (d)~verify that the
        dual solution taken from the tableau is dual feasible with
        $W=Z^{*}$.
  \item (Price theorem.) In the LP of Problem~14 the stock of resource A
        rises from $b_{1}=4$ to $b_{1}=5$. Predict the new optimal value
        with the price theorem, then verify the prediction by re-solving
        the LP (the same two constraints remain binding).
  \item Purchasing decisions with the shadow prices $y_{1}^{*}=1$,
        $y_{2}^{*}=1$ from Problem~14 (all prices in units of $Z$):
        (a)~a supplier offers one extra unit of resource A at price
        $1.25$ (accept?
        (b)~one extra unit of resource B at price $0.8$) accept?
        (c)~In the LP of Problem~10, extra units of resource 3 are offered
        at price $0.1$ each: accept?
  \item A joiner makes tables ($x_{1}$) and shelves ($x_{2}$). A table
        earns profit $3$ and needs $1$ unit of wood and $2$ machine-hours;
        a shelf earns $2$ and needs $1$ unit of wood and $1$ machine-hour.
        Available: $4$ units of wood, $6$ machine-hours: this is exactly
        the LP of Problem~14 with $Z^{*}=10$ and $y^{*}=(1,1)$.
        (a)~Write the dual.
        (b)~Interpret the dual variables, the dual constraints, and the
        dual objective economically.
        (c)~What is the smallest total payment for which the joiner should
        be willing to rent out his entire stock of wood and machine time?
  \item In Problem~10 the third constraint has slack $2$ at the optimum
        and $y_{3}^{*}=0$. Explain this value twice: once through
        complementary slackness and once economically. Does
        $y_{i}^{*}=0$ conversely always imply that constraint $i$ has
        slack?
\end{problems}

\section{Solutions}

\begin{solutions}
  \item Two constraints give two dual variables $y_{1},y_{2}\ge 0$; the
        right-hand sides $(4,6)$ become the dual objective; the columns of
        $x_{1}$ and $x_{2}$, namely $(1,2)$ and $(1,1)$, become the dual
        constraint rows with right-hand sides $c=(3,2)$:
        \[
          \min\; W=4y_{1}+6y_{2}
          \quad\text{s.t.}\quad
          y_{1}+2y_{2}\ge 3,\qquad y_{1}+y_{2}\ge 2,\qquad
          y_{1},y_{2}\ge 0 .
        \]

  \item Three constraints, two variables, so the dual has three variables
        and two constraints:
        \[
          \min\; W=10y_{1}+6y_{2}+10y_{3}
          \quad\text{s.t.}\quad
          2y_{1}+y_{2}+y_{3}\ge 4,\qquad
          y_{1}+y_{2}+2y_{3}\ge 3,\qquad
          y_{1},y_{2},y_{3}\ge 0 .
        \]

  \item A minimum problem with $\ge$ constraints and non-negative
        variables dualises to a maximum problem with $\le$ constraints and
        non-negative variables. One dual variable per constraint,
        objective from the right-hand sides $(3,2,2)$, constraints from
        the variable columns $(3,1,1)$ and $(1,1,2)$ with right-hand
        sides $(6,4)$:
        \[
          \max\; W=3u_{1}+2u_{2}+2u_{3}
          \quad\text{s.t.}\quad
          3u_{1}+u_{2}+u_{3}\le 6,\qquad
          u_{1}+u_{2}+2u_{3}\le 4,\qquad
          u_{1},u_{2},u_{3}\ge 0 .
        \]

  \item Sign rules for a max primal: the $\le$ constraint gives
        $y_{1}\ge 0$, the $\ge$ constraint gives $y_{2}\le 0$, the
        equality gives $y_{3}$ free. Both variables satisfy
        $x_{j}\ge 0$, so both dual constraints are of type $\ge$:
        \[
          \min\; W=12y_{1}+2y_{2}+8y_{3}
          \quad\text{s.t.}\quad
          y_{1}+y_{2}+y_{3}\ge 2,\qquad
          3y_{1}-y_{2}+y_{3}\ge 5,
        \]
        with $y_{1}\ge 0$, $y_{2}\le 0$, $y_{3}$ free.

  \item Both constraints are $\le$, so $y_{1},y_{2}\ge 0$. The variable
        $x_{1}\ge 0$ gives a $\ge$ dual constraint, while the \emph{free}
        variable $x_{2}$ tightens its dual constraint to an equality:
        \[
          \min\; W=8y_{1}+10y_{2}
          \quad\text{s.t.}\quad
          2y_{1}+y_{2}\ge 1,\qquad
          y_{1}+2y_{2}= 4,\qquad
          y_{1},y_{2}\ge 0 .
        \]

  \item The dual of the given LP is
        \[
          \min\; W=24y_{1}+6y_{2}
          \quad\text{s.t.}\quad
          6y_{1}+y_{2}\ge 5,\qquad 4y_{1}+2y_{2}\ge 4,\qquad
          y_{1},y_{2}\ge 0 .
        \]
        Now dualise this minimum problem: its two $\ge$ constraints give
        variables $u_{1},u_{2}\ge 0$; its right-hand sides $(5,4)$ become
        the new objective; its variables $y_{1},y_{2}\ge 0$ give $\le$
        constraints with right-hand sides $(24,6)$, rows taken from the
        columns of $y_{1}$ and $y_{2}$:
        \[
          \max\; 5u_{1}+4u_{2}
          \quad\text{s.t.}\quad
          6u_{1}+4u_{2}\le 24,\qquad u_{1}+2u_{2}\le 6,\qquad
          u_{1},u_{2}\ge 0 .
        \]
        Renaming $u$ as $x$ recovers the primal exactly.

  \item Primal feasibility of $\hat{x}=(1,2)$: $1+2=3\le 4$ and
        $2\cdot 1+2=4\le 6$; also $\hat{x}\ge 0$. Dual feasibility of
        $\hat{y}=(1,1)$ (dual from Solution~1): $1+2\cdot 1=3\ge 3$ and
        $1+1=2\ge 2$; $\hat{y}\ge 0$. Objective values:
        \[
          Z(\hat{x})=3\cdot 1+2\cdot 2=7,
          \qquad
          W(\hat{y})=4\cdot 1+6\cdot 1=10 .
        \]
        Indeed $7\le 10$, as weak duality requires. Bounds: every
        dual-feasible point caps the primal optimum and every
        primal-feasible point supports it, so
        $7\le Z^{*}\le 10$.

  \item Feasibility of $x=(2,2)$: $2+2=4\le 4$ (binds) and
        $2\cdot 2+2=6\le 6$ (binds); $Z=3\cdot 2+2\cdot 2=10$.
        Feasibility of $y=(1,1)$: checked in Solution~7; $W=10$.
        Since $Z(x)=10=W(y)$, weak duality closes the gap: for any
        primal-feasible $x'$ we have $Z(x')\le W(y)=10=Z(x)$, so $x$ is
        primal-optimal; for any dual-feasible $y'$ we have
        $W(y')\ge Z(x)=10=W(y)$, so $y$ is dual-optimal. Hence
        $Z^{*}=W^{*}=10$.

  \item (a)~The points $(t,t)$, $t\ge 0$, are feasible since
        $t-t=0\le 1$, and $Z=2t\to\infty$: the LP is unbounded.
        (b)~The dual has one variable $y_{1}\ge 0$ and one constraint per
        primal variable:
        \[
          \min\; W=y_{1}
          \quad\text{s.t.}\quad
          y_{1}\ge 1,\qquad -y_{1}\ge 1,\qquad y_{1}\ge 0 .
        \]
        The second constraint forces $y_{1}\le -1$, contradicting
        $y_{1}\ge 0$: the dual is infeasible.
        (c)~If the dual had any feasible point $y$, weak duality would
        give $Z(x)\le W(y)<\infty$ for all primal-feasible $x$, so $Z$
        could not grow without bound. An unbounded primal therefore
        forces an infeasible dual.

  \item First classify the constraints at $x^{*}=(4,2)$:
        \[
          2\cdot 4+2=10 \;(\text{binds}),\qquad
          4+2=6 \;(\text{binds}),\qquad
          4+2\cdot 2=8<10 \;(\text{slack }2).
        \]
        Complementary slackness: the slack third constraint forces
        $y_{3}=0$; the positive variables $x_{1}^{*}=4>0$ and
        $x_{2}^{*}=2>0$ force both dual constraints (Solution~2) to bind:
        \[
          2y_{1}+y_{2}+y_{3}=4,\qquad y_{1}+y_{2}+2y_{3}=3 .
        \]
        With $y_{3}=0$: $2y_{1}+y_{2}=4$ and $y_{1}+y_{2}=3$; subtracting
        gives $y_{1}=1$, hence $y_{2}=2$. So $y^{*}=(1,2,0)$, which is
        non-negative and dual feasible. Strong duality:
        \[
          W^{*}=10\cdot 1+6\cdot 2+10\cdot 0=22
          = 4\cdot 4+3\cdot 2 = Z^{*}. \checkmark
        \]

  \item The dual is
        $\min\; W=4y_{1}+6y_{2}$ s.t.\ $y_{1}+y_{2}\ge 2$,
        $y_{1}+3y_{2}\ge 1$, $y\ge 0$. At $y^{*}=(2,0)$ the first dual
        constraint binds ($2+0=2$) and the second is strict
        ($2+0=2>1$). Complementary slackness backwards:
        \begin{itemize}[leftmargin=1.6em,itemsep=0.3ex]
          \item $y_{1}^{*}=2>0$ $\Rightarrow$ constraint 1 binds:
                $x_{1}+x_{2}=4$;
          \item dual constraint 2 strict $\Rightarrow$ $x_{2}^{*}=0$.
        \end{itemize}
        Hence $x_{1}^{*}=4$, $x^{*}=(4,0)$. Feasibility check:
        $4\le 4$ (binds) and $4+0=4\le 6$ (slack, consistent with
        $y_{2}^{*}=0$). Values:
        $Z^{*}=2\cdot 4+1\cdot 0=8$ and $W^{*}=4\cdot 2+6\cdot 0=8$.
        $\checkmark$

  \item Feasibility: $x=(3,0)$ satisfies $3\le 4$ (slack $1$) and
        $6\le 6$ (binds), with $Z=9$; $y=(1,1)$ is dual feasible with
        $W=10$ (Solution~7). Now the four conditions:
        \[
          y_{1}(4-x_{1}-x_{2})=1\cdot(4-3)=1\ne 0
          \quad\textbf{(violated)},
        \]
        \[
          y_{2}(6-2x_{1}-x_{2})=1\cdot 0=0,\qquad
          x_{1}(y_{1}+2y_{2}-3)=3\cdot 0=0,\qquad
          x_{2}(y_{1}+y_{2}-2)=0\cdot 0=0 .
        \]
        One condition fails, so the pair is \emph{not} jointly optimal.
        Consistently, the duality gap is open: $Z=9<10=W$. (In fact $y$
        is dual-optimal by Solution~8, so it is $x$ that is not optimal.)

  \item Dual (min primal with $\ge$ constraints $\Rightarrow$ max dual
        with $\le$ constraints):
        \[
          \max\; W=5u_{1}+6u_{2}
          \quad\text{s.t.}\quad
          u_{1}+u_{2}\le 2,\qquad u_{1}+2u_{2}\le 3,\qquad
          u_{1},u_{2}\ge 0 .
        \]
        At $x^{*}=(4,1)$ both primal constraints bind:
        $4+1=5$ and $4+2\cdot 1=6$. Since $x_{1}^{*}=4>0$ and
        $x_{2}^{*}=1>0$, both dual constraints must bind:
        \[
          u_{1}+u_{2}=2,\qquad u_{1}+2u_{2}=3
          \quad\Longrightarrow\quad u_{2}=1,\; u_{1}=1 .
        \]
        So $u^{*}=(1,1)\ge 0$, dual feasible, and
        \[
          W^{*}=5\cdot 1+6\cdot 1=11
          =2\cdot 4+3\cdot 1=Z^{*}. \checkmark
        \]

  \item (a)~Basic variables and their RHS values: $x_{1}=2$, $x_{2}=2$;
        the slacks are non-basic, hence zero. Optimal solution
        $x^{*}=(2,2)$ with $Z^{*}=10$ (bottom right).
        (b)~$s_{1}=s_{2}=0$: both constraints bind.
        (c)~Shadow prices are the bottom-row entries under the slack
        columns: $y_{1}^{*}=1$, $y_{2}^{*}=1$.
        (d)~Dual (Solution~1): $y_{1}+2y_{2}=1+2=3\ge 3$ and
        $y_{1}+y_{2}=2\ge 2$, both satisfied (as equalities), $y\ge 0$;
        and $W=4\cdot 1+6\cdot 1=10=Z^{*}$. $\checkmark$

  \item Price theorem: $\Delta Z^{*}=y_{1}^{*}\cdot\Delta b_{1}
        =1\cdot 1=1$, so the predicted optimum is $Z^{*}=10+1=11$.
        Verification: with the same binding set,
        \[
          x_{1}+x_{2}=5,\qquad 2x_{1}+x_{2}=6
          \quad\Longrightarrow\quad x_{1}=1,\; x_{2}=4 ,
        \]
        both non-negative, and $Z=3\cdot 1+2\cdot 4=11$. $\checkmark$
        The prediction is exact here because the optimal basis is
        unchanged; for large enough increases of $b_{1}$ the basis would
        switch and the rate $y_{1}^{*}$ would no longer apply.

  \item (a)~No: the price $1.25$ exceeds the shadow price
        $y_{1}^{*}=1$; the extra unit would raise $Z$ by $1$ but cost
        $1.25$, a net loss of $0.25$.
        (b)~Yes: $0.8<y_{2}^{*}=1$; net gain $1-0.8=0.2$ per unit.
        (c)~No: resource 3 has shadow price $y_{3}^{*}=0$ (its constraint
        already has slack), so an extra unit adds nothing to $Z^{*}$;
        any positive price is too high.

  \item (a)~Dual:
        \[
          \min\; W=4y_{1}+6y_{2}
          \quad\text{s.t.}\quad
          y_{1}+2y_{2}\ge 3,\qquad y_{1}+y_{2}\ge 2,\qquad
          y_{1},y_{2}\ge 0 .
        \]
        (b)~$y_{1}$ is the rental price offered per unit of wood, $y_{2}$
        per machine-hour. The first dual constraint prices the resource
        bundle of one table ($1$ wood, $2$ machine-hours) at no less than
        the table's profit $3$; the second does the same for a shelf. If
        such a constraint were violated, the joiner would refuse the
        offer and produce that good instead. The dual objective
        $4y_{1}+6y_{2}$ is the bidder's total rent for the whole stock,
        which the bidder minimises subject to making an acceptable offer.
        (c)~By strong duality the cheapest acceptable offer equals the
        maximal profit: $W^{*}=Z^{*}=10$. Indeed at $y^{*}=(1,1)$:
        $4\cdot 1+6\cdot 1=10$. Below a total of $10$ the joiner does
        better producing himself.

  \item Complementary slackness: the condition
        $y_{3}^{*}\,(b_{3}-a_{3}^{\top}x^{*})=0$ with slack
        $b_{3}-a_{3}^{\top}x^{*}=10-8=2>0$ forces $y_{3}^{*}=0$.
        Economically: two units of resource 3 are already lying idle at
        the optimum, so an additional unit cannot raise profit: its
        marginal value, the shadow price, is zero. The converse is
        \emph{not} valid in general: in a degenerate optimum a
        \emph{binding} constraint can also have shadow price zero. Only
        the direction ``slack $\Rightarrow$ zero price'' is guaranteed.
\end{solutions}

\section{Additional Problems}

\begin{problems}
  \item Write the dual of
        \[
          \max\; Z=6x_{1}+x_{2}+2x_{3}
          \quad\text{s.t.}\quad
          x_{1}+2x_{2}+x_{3}\le 5,\qquad
          2x_{1}+x_{2}+3x_{3}\le 8,\qquad
          x_{1},x_{2},x_{3}\ge 0 .
        \]
  \item Using the general sign rules, write the dual of
        \[
          \min\; Z=4x_{1}+2x_{2}
          \quad\text{s.t.}\quad
          x_{1}+x_{2}\ge 3,\qquad
          x_{1}-x_{2}\le 1,\qquad
          x_{1}+2x_{2}= 6,\qquad x_{1},x_{2}\ge 0 .
        \]
  \item For
        $\max\; Z=2x_{1}+3x_{2}$ s.t.\ $x_{1}+2x_{2}\le 6$,
        $x_{1}+x_{2}\le 4$, $x_{1},x_{2}\ge 0$:
        verify that $\hat{x}=(2,1)$ is primal feasible and
        $\hat{y}=(1,1)$ is dual feasible, state the resulting bounds on
        $Z^{*}$, and then show that $x=(2,2)$ closes the gap, proving
        $\hat{y}$ optimal.
  \item The LP of Additional Problem~3 has the optimal solution
        $x^{*}=(2,2)$. Use complementary slackness to compute $y^{*}$ and
        verify strong duality.
  \item Consider
        $\max\; Z=x_{1}+2x_{2}$ s.t.\ $x_{1}+x_{2}\le 3$, $x_{2}\le 2$,
        $x_{1},x_{2}\ge 0$. Its dual has the optimal solution
        $y^{*}=(1,1)$. Recover the optimal primal solution via
        complementary slackness and verify the objective values.
  \item Solving
        $\max\; Z=2x_{1}+3x_{2}$ s.t.\
        $x_{1}+2x_{2}\le 6$ (wood, slack $s_{1}$),
        $x_{1}+x_{2}\le 4$ (labour, slack $s_{2}$),
        $x_{1},x_{2}\ge 0$
        yields the end tableau
        \begin{center}
        \begin{tabular}{c|cccc|c}
          Basis & $x_{1}$ & $x_{2}$ & $s_{1}$ & $s_{2}$ & RHS\\
          \hline
          $x_{1}$ & $1$ & $0$ & $-1$ & $2$  & $2$\\
          $x_{2}$ & $0$ & $1$ & $1$  & $-1$ & $2$\\
          \hline
          $Z$     & $0$ & $0$ & $1$  & $1$  & $10$\\
        \end{tabular}
        \end{center}
        Read off the optimal solution, the binding constraints, and the
        shadow prices, and verify dual feasibility and $W=Z^{*}$.
  \item (Price theorem.) In the LP of Additional Problem~6 the labour
        stock rises from $b_{2}=4$ to $b_{2}=5$. Predict the new optimal
        value with the price theorem and verify by re-solving (the same
        constraints remain binding).
  \item Economics of the LP of Additional Problem~6 ($6$ wood, $4$
        labour, $y^{*}=(1,1)$, $Z^{*}=10$):
        (a)~What is the smallest total payment for which the firm should
        hand over its entire resource stock?
        (b)~An extra unit of wood is offered at price $0.5$: accept, and
        with what net gain?
        (c)~What does it mean if a resource is offered at a price exactly
        equal to its shadow price?
  \item A firm makes three products from machine hours and labour hours:
        \[
          \max\; Z=11x_{1}+14x_{2}+13x_{3}
          \quad\text{s.t.}\quad
          3x_{1}+2x_{2}+4x_{3}\le 60,\qquad
          x_{1}+4x_{2}+x_{3}\le 60,\qquad x\ge 0 .
        \]
        (a)~Write the dual and solve it graphically.
        (b)~Recover the primal optimum by complementary slackness and
        verify $Z^{*}=W^{*}$.
        (c)~Which products are made, which resources are left over, and
        which margin would product~3 need before it is worth producing?
        (d)~State the three dual slacks and interpret them.
  \item For
        \[
          \max\; Z=4x_{1}+6x_{2}+5x_{3}
          \quad\text{s.t.}\quad
          \begin{array}{rcll}
            x_{1}+2x_{2}+x_{3} &\le& 20 & \text{(resource A)}\\
            3x_{1}+x_{2}+2x_{3} &\le& 90 & \text{(resource B)}\\
            x &\ge& 0,
          \end{array}
        \]
        solve the dual graphically, recover the primal optimum by
        complementary slackness, and answer: which products are made, how
        much of each resource is left over, and by how much would the
        margin of each product not made have to rise?
  \item A feed is mixed from two grains at minimum cost:
        \[
          \min\; C=3x_{1}+2x_{2}
          \quad\text{s.t.}\quad
          x_{1}+x_{2}\ge 4,\qquad x_{1}+3x_{2}\ge 4,\qquad
          2x_{1}+x_{2}\ge 5,\qquad x\ge 0 .
        \]
        (a)~Write the dual and explain why it is a standard maximum
        problem. (b)~Solve the dual by the simplex method (ties: leftmost
        column). (c)~Read the optimal mix, the minimal cost, and the
        surplus of each requirement off the end tableau. (d)~Confirm the
        result by complementary slackness.
  \item A firm solves
        \[
          \max\; Z=5x_{1}+4x_{2}
          \quad\text{s.t.}\quad
          \begin{array}{rcll}
            x_{1}+x_{2} &\le& 10 & \text{(labour)}\\
            2x_{1}+x_{2} &\le& 16 & \text{(machine)}\\
            x_{2} &\ge& 6 & \text{(contract)}\\
            x &\ge& 0,
          \end{array}
        \]
        and finds the optimum $x^{*}=(4,6)$. (a)~Write the dual with the
        general sign rules. (b)~Find $y^{*}$ by complementary slackness;
        check the signs, dual feasibility and $W^{*}=Z^{*}$.
        (c)~Interpret the dual value of the contract row, and confirm it by
        re-solving with a contract of $7$ units. (d)~Which price does the
        contract row carry if it is first written as $-x_{2}\le -6$?
  \item The minimum problem
        \[
          \min\; C=2x_{1}+5x_{2}
          \quad\text{s.t.}\quad
          x_{1}+2x_{2}\ge 8,\qquad x_{1}\le 4,\qquad
          x_{1}-3x_{2}\ge -10,\qquad x\ge 0
        \]
        has the optimum $x^{*}=(4,2)$. (a)~Write the dual with the
        general sign rules. (b)~Find $u^{*}$ by complementary slackness
        and check signs, feasibility and $W^{*}=C^{*}$. (c)~Interpret
        $u_{1}^{*}$, $u_{2}^{*}$ and $u_{3}^{*}$, and confirm $u_{2}^{*}$
        by re-solving with the capacity $x_{1}\le 5$.
  \item (Kuhn--Tucker.) For the LP of Practice Problem~14,
        $\max Z=3x_{1}+2x_{2}$ s.t.\ $x_{1}+x_{2}\le 4$,
        $2x_{1}+x_{2}\le 6$, $x\ge 0$:
        (a)~write the Lagrangian with multipliers $\lambda_{1},\lambda_{2}$;
        (b)~list the Kuhn--Tucker conditions (with $x\ge 0$);
        (c)~check them at $x=(2,2)$, $\lambda=(1,1)$;
        (d)~explain which of the conditions are the dual constraints,
        which are complementary slackness, and why here they are also
        sufficient for optimality.
  \item (Kuhn--Tucker with explicit sign multipliers.) For the LP of
        Additional Problem~9 write the Lagrangian with multipliers
        $\lambda_{1},\lambda_{2}$ for the resources and $\mu_{1},\mu_{2},
        \mu_{3}$ for $x_{j}\ge 0$. Show that $x=(12,12,0)$,
        $\lambda=(3,2)$, $\mu=(0,0,1)$ satisfy every Kuhn--Tucker condition,
        and interpret $\mu_{3}$.
  \item The LP of Additional Problem~10 has the dual slacks
        $t=(1,4,0)$ at its optimum. True or false (justify):
        (a)~raising the margin of product~2 from $6$ to $9$ makes it worth
        producing;
        (b)~raising the margin of product~1 from $4$ to $5.5$ makes it
        worth producing;
        (c)~product~3 has dual slack $0$, so its margin may fall by any
        amount without changing the plan;
        (d)~at the optimum each dual slack equals the absolute value of
        the reduced cost of its variable.
\end{problems}

\section{Additional Solutions}

\begin{solutions}
  \item Two constraints give $y_{1},y_{2}\ge 0$; three variables give
        three $\ge$ constraints with right-hand sides $(6,1,2)$:
        \[
          \min\; W=5y_{1}+8y_{2}
          \quad\text{s.t.}\quad
          y_{1}+2y_{2}\ge 6,\qquad
          2y_{1}+y_{2}\ge 1,\qquad
          y_{1}+3y_{2}\ge 2,\qquad
          y_{1},y_{2}\ge 0 .
        \]

  \item For a min primal the dual is a max problem. Sign rules: the
        $\ge$ constraint gives $u_{1}\ge 0$, the $\le$ constraint gives
        $u_{2}\le 0$, the equality gives $u_{3}$ free; the variables
        $x_{1},x_{2}\ge 0$ give $\le$ dual constraints:
        \[
          \max\; W=3u_{1}+u_{2}+6u_{3}
          \quad\text{s.t.}\quad
          u_{1}+u_{2}+u_{3}\le 4,\qquad
          u_{1}-u_{2}+2u_{3}\le 2,
        \]
        with $u_{1}\ge 0$, $u_{2}\le 0$, $u_{3}$ free.

  \item Primal feasibility of $\hat{x}=(2,1)$: $2+2=4\le 6$,
        $2+1=3\le 4$; $Z(\hat{x})=2\cdot 2+3\cdot 1=7$. The dual is
        \[
          \min\; W=6y_{1}+4y_{2}
          \quad\text{s.t.}\quad
          y_{1}+y_{2}\ge 2,\qquad 2y_{1}+y_{2}\ge 3,\qquad y\ge 0,
        \]
        and $\hat{y}=(1,1)$ satisfies $1+1=2\ge 2$, $2+1=3\ge 3$;
        $W(\hat{y})=6+4=10$. Weak duality gives $7\le Z^{*}\le 10$.
        The point $x=(2,2)$ is feasible ($2+4=6\le 6$, $2+2=4\le 4$)
        with $Z=2\cdot 2+3\cdot 2=10=W(\hat{y})$: the certificate
        argument makes $x$ primal-optimal and $\hat{y}$ dual-optimal,
        so $Z^{*}=W^{*}=10$.

  \item At $x^{*}=(2,2)$ both constraints bind ($2+4=6$, $2+2=4$), and
        both variables are positive, so both dual constraints bind:
        \[
          y_{1}+y_{2}=2,\qquad 2y_{1}+y_{2}=3
          \quad\Longrightarrow\quad y_{1}=1,\; y_{2}=1 .
        \]
        $y^{*}=(1,1)\ge 0$ is dual feasible, and
        $W^{*}=6\cdot 1+4\cdot 1=10=2\cdot 2+3\cdot 2=Z^{*}$.
        $\checkmark$

  \item Dual:
        $\min\; W=3y_{1}+2y_{2}$ s.t.\ $y_{1}\ge 1$,
        $y_{1}+y_{2}\ge 2$, $y\ge 0$. At $y^{*}=(1,1)$:
        $y_{1}^{*}=1>0$ forces $x_{1}+x_{2}=3$, and $y_{2}^{*}=1>0$
        forces $x_{2}=2$. Hence $x_{1}=1$, i.e.\ $x^{*}=(1,2)$, which is
        feasible. Values:
        $Z^{*}=1+2\cdot 2=5$ and $W^{*}=3\cdot 1+2\cdot 1=5$.
        $\checkmark$

  \item Basic variables: $x_{1}=2$, $x_{2}=2$, so $x^{*}=(2,2)$ and
        $Z^{*}=10$; the slacks are non-basic and zero, so both
        constraints bind. Shadow prices from the bottom row under the
        slack columns: $y_{1}^{*}=1$ (wood), $y_{2}^{*}=1$ (labour).
        Dual feasibility (dual from Additional Solution~3):
        $y_{1}+y_{2}=2\ge 2$ and $2y_{1}+y_{2}=3\ge 3$, both binding,
        $y\ge 0$; and $W=6\cdot 1+4\cdot 1=10=Z^{*}$. $\checkmark$

  \item Prediction: $\Delta Z^{*}=y_{2}^{*}\cdot\Delta b_{2}=1\cdot 1=1$,
        so $Z^{*}=10+1=11$. Verification with the same binding set:
        \[
          x_{1}+2x_{2}=6,\qquad x_{1}+x_{2}=5
          \quad\Longrightarrow\quad x_{2}=1,\; x_{1}=4 ,
        \]
        both non-negative, and $Z=2\cdot 4+3\cdot 1=11$. $\checkmark$

  \item (a)~By strong duality the resource stock is worth exactly
        $W^{*}=6y_{1}^{*}+4y_{2}^{*}=10$: any total payment below $10$
        is worse than producing, any payment of at least $10$ is
        acceptable.
        (b)~Accept: the shadow price of wood is $1$, so the extra unit
        raises $Z^{*}$ by $1$ at a cost of $0.5$, net gain $0.5$.
        (c)~At a price equal to the shadow price the firm is exactly
        indifferent at the margin: the extra unit raises profit by
        precisely what it costs.

  \item (a)~Two constraints give two dual variables, three products give
        three dual constraints:
        \[
          \min\; W=60y_{1}+60y_{2}
          \quad\text{s.t.}\quad
          3y_{1}+y_{2}\ge 11,\qquad 2y_{1}+4y_{2}\ge 14,\qquad
          4y_{1}+y_{2}\ge 13,\qquad y\ge 0 .
        \]
        In the $(y_{1},y_{2})$-plane the feasible set lies above the three
        lines $y_{2}=11-3y_{1}$, $y_{2}=(7-y_{1})/2$ and $y_{2}=13-4y_{1}$.
        Its corners: on the $y_{2}$-axis the binding line is the third
        ($y_{2}\ge 13$), giving $(0,13)$; the third and first lines meet at
        $4y_{1}+y_{2}=13$, $3y_{1}+y_{2}=11$, i.e.\ $(2,5)$ (and
        $2\cdot2+4\cdot5=24\ge 14$); the first and second meet at
        $3y_{1}+y_{2}=11$, $y_{1}+2y_{2}=7$, i.e.\ $(3,2)$ (and
        $4\cdot3+2=14\ge 13$); the second meets the $y_{1}$-axis at
        $(7,0)$. Values of $W=60(y_{1}+y_{2})$:
        \[
          W(0,13)=780,\quad W(2,5)=420,\quad W(3,2)=300,\quad W(7,0)=420 .
        \]
        The region is unbounded upwards, but $W$ grows there, so the
        minimum sits at the corner $y^{*}=(3,2)$ with $W^{*}=300$.
        \begin{gallery}
        \hfill
        \begin{tikzpicture}
          \begin{axis}[geoaxis, xmin=-0.6, xmax=8.6, ymin=-0.8, ymax=14.6,
              xtick={2,4,6,8}, ytick={2,4,6,8,10,12,14},
              xlabel={$y_{1}$}, ylabel={$y_{2}$},
              title={dual feasible set}]
            \addplot[draw=none, name path=lA, domain=0:2, samples=2]{13-4*x};
            \addplot[draw=none, name path=uA, domain=0:2, samples=2]{14.4};
            \region{lA}{uA}{}
            \addplot[draw=none, name path=lB, domain=2:3, samples=2]{11-3*x};
            \addplot[draw=none, name path=uB, domain=2:3, samples=2]{14.4};
            \region{lB}{uB}{}
            \addplot[draw=none, name path=lC, domain=3:7, samples=2]{(7-x)/2};
            \addplot[draw=none, name path=uC, domain=3:7, samples=2]{14.4};
            \region{lC}{uC}{}
            \addplot[draw=none, name path=lD, domain=7:8.4, samples=2]{0};
            \addplot[draw=none, name path=uD, domain=7:8.4, samples=2]{14.4};
            \region{lD}{uD}{}
            \addplot[domain=0:4.2, samples=2]{11-3*x};
            \addplot[domain=0:8.4, samples=2]{(7-x)/2};
            \addplot[domain=0:3.5, samples=2]{13-4*x};
            \addplot[gray, dashed, thin, domain=0:5, samples=2]{5-x};
            \addplot[only marks, mark=*, mark size=1.4pt, black, forget plot]
              coordinates {(0,13) (2,5) (3,2) (7,0)};
            \node[font=\tiny, anchor=south west] at (axis cs:3.1,2.1) {$y^{*}=(3,2)$};
            \node[font=\tiny, anchor=west] at (axis cs:2.15,5.2) {$(2,5)$};
          \end{axis}
        \end{tikzpicture}
        \hfill\null
        \gcap{The dual feasible set (shaded, unbounded upwards) and the
        lowest iso-cost line $y_{1}+y_{2}=5$ (dashed) through $y^{*}=(3,2)$.}
        \end{gallery}

        (b)~At $y^{*}=(3,2)$ the third dual constraint is slack
        ($4\cdot3+2=14>13$), so $x_{3}^{*}=0$. Both $y_{i}^{*}>0$, so both
        primal rows bind:
        \[
          3x_{1}+2x_{2}=60,\qquad x_{1}+4x_{2}=60 .
        \]
        From the second, $x_{1}=60-4x_{2}$; substituting,
        $180-12x_{2}+2x_{2}=60$, so $x_{2}=12$ and $x_{1}=12$. Both are
        nonnegative, and
        $Z^{*}=11\cdot12+14\cdot12=132+168=300=W^{*}$. $\checkmark$

        (c)~Products $1$ and $2$ are made, $12$ units each; product~3 is
        not made. Both resources are used up completely ($36+24=60$,
        $12+48=60$), which fits their positive shadow prices: nothing is
        left over. One unit of product~3 uses $4$ machine hours and $1$
        labour hour, worth $4\cdot3+1\cdot2=14$ at shadow prices; its
        margin $13$ must rise above $14$ before it is worth producing.

        (d)~$t_{j}=(A^{\top}y^{*})_{j}-c_{j}$:
        $t_{1}=11-11=0$, $t_{2}=14-14=0$, $t_{3}=14-13=1$. Products $1$
        and $2$ break even exactly at shadow prices (and are made);
        product~3 falls short by $1$ per unit, which is the margin
        increase it needs.

  \item The dual is
        \[
          \min\; W=20y_{1}+90y_{2}
          \quad\text{s.t.}\quad
          y_{1}+3y_{2}\ge 4,\qquad 2y_{1}+y_{2}\ge 6,\qquad
          y_{1}+2y_{2}\ge 5,\qquad y\ge 0 .
        \]
        Corners of the feasible set: on the $y_{2}$-axis the binding line
        is the second ($y_{2}\ge 6$), giving $(0,6)$; the second and third
        lines meet at $2y_{1}+y_{2}=6$, $y_{1}+2y_{2}=5$, i.e.\
        $(\tfrac73,\tfrac43)$ (and $\tfrac73+4=\tfrac{19}{3}\ge 4$); on the
        $y_{1}$-axis the binding line is the third ($y_{1}\ge 5$), giving
        $(5,0)$. (The first line is cut off everywhere: for example its
        crossing with the second line, $(\tfrac{14}{5},\tfrac25)$, violates
        the third.) Values:
        \[
          W(0,6)=540,\qquad W(\tfrac73,\tfrac43)=\tfrac{140}{3}+120
          =\tfrac{500}{3},\qquad W(5,0)=100 .
        \]
        So $y^{*}=(5,0)$, $W^{*}=100$.

        Complementary slackness: at $(5,0)$ the first dual constraint is
        slack ($5>4$, $t_{1}=1$) and so is the second ($10>6$, $t_{2}=4$),
        hence $x_{1}^{*}=x_{2}^{*}=0$. Since $y_{1}^{*}=5>0$, resource~A
        binds: $x_{3}=20$. Since $y_{2}^{*}=0$, resource~B may have slack:
        $2\cdot20=40\le 90$, slack $50$. So $x^{*}=(0,0,20)$ and
        $Z^{*}=5\cdot20=100=W^{*}$. $\checkmark$

        Only product~3 is made ($20$ units). Resource~A is used up;
        $50$ units of resource~B are left over, which is why B's shadow
        price is $0$. Product~1 would need a margin above
        $1\cdot5=5$, i.e.\ $1$ more; product~2 above $2\cdot5=10$, i.e.\
        $4$ more: exactly the dual slacks.

  \item (a)~A min problem with $\ge$ rows and $x\ge 0$ has a max dual with
        $\le$ rows and $u\ge 0$:
        \[
          \max\; W=4u_{1}+4u_{2}+5u_{3}
          \quad\text{s.t.}\quad
          u_{1}+u_{2}+2u_{3}\le 3,\qquad u_{1}+3u_{2}+u_{3}\le 2,\qquad
          u\ge 0 .
        \]
        Its right-hand sides are the costs $3$ and $2$, both
        nonnegative, so it is a standard maximum problem: the origin is
        feasible and the slacks $t_{1},t_{2}$ form the first basis.

        (b)~Initial tableau:
        \begin{center}
        \begin{tabular}{c|ccccc|c}
        Basis & $u_{1}$ & $u_{2}$ & $u_{3}$ & $t_{1}$ & $t_{2}$ & RHS\\
        \hline
        $t_{1}$ & $1$ & $1$ & $2$ & $1$ & $0$ & $3$\\
        $t_{2}$ & $1$ & $3$ & $1$ & $0$ & $1$ & $2$\\
        \hline
        $W$ & $-4$ & $-4$ & $-5$ & $0$ & $0$ & $0$\\
        \end{tabular}
        \end{center}
        \emph{Iteration 1.} Pivot column $u_{3}$ ($-5$). Ratios:
        $3/2=1.5$, $2/1=2$; $t_{1}$ leaves, pivot element $2$. Row
        operations: $R_{1}\to R_{1}/2$; then $R_{2}\to R_{2}-R_{1}$,
        $R_{W}\to R_{W}+5R_{1}$.
        \begin{center}
        \begin{tabular}{c|ccccc|c}
        Basis & $u_{1}$ & $u_{2}$ & $u_{3}$ & $t_{1}$ & $t_{2}$ & RHS\\
        \hline
        $u_{3}$ & $\tfrac12$ & $\tfrac12$ & $1$ & $\tfrac12$ & $0$ & $\tfrac32$\\
        $t_{2}$ & $\tfrac12$ & $\tfrac52$ & $0$ & $-\tfrac12$ & $1$ & $\tfrac12$\\
        \hline
        $W$ & $-\tfrac32$ & $-\tfrac32$ & $0$ & $\tfrac52$ & $0$ & $\tfrac{15}{2}$\\
        \end{tabular}
        \end{center}
        \emph{Iteration 2.} Tie between $u_{1}$ and $u_{2}$; take $u_{1}$.
        Ratios: $R_{1}$: $\tfrac32\big/\tfrac12=3$; $R_{2}$:
        $\tfrac12\big/\tfrac12=1$; $t_{2}$ leaves, pivot element
        $\tfrac12$. Row operations: $R_{2}\to 2R_{2}$; then
        $R_{1}\to R_{1}-\tfrac12R_{2}$, $R_{W}\to R_{W}+\tfrac32R_{2}$.
        \begin{center}
        \begin{tabular}{c|ccccc|c}
        Basis & $u_{1}$ & $u_{2}$ & $u_{3}$ & $t_{1}$ & $t_{2}$ & RHS\\
        \hline
        $u_{3}$ & $0$ & $-2$ & $1$ & $1$ & $-1$ & $1$\\
        $u_{1}$ & $1$ & $5$ & $0$ & $-1$ & $2$ & $1$\\
        \hline
        $W$ & $0$ & $6$ & $0$ & $1$ & $3$ & $9$\\
        \end{tabular}
        \end{center}
        No negative entry: optimal, $u^{*}=(1,0,1)$, $W^{*}=9$.

        (c)~The primal optimum stands in the $W$-row under the dual's slack
        columns: $x_{1}^{*}=1$ (under $t_{1}$), $x_{2}^{*}=3$ (under
        $t_{2}$), with minimal cost $C^{*}=W^{*}=9$; check
        $3\cdot1+2\cdot3=9$. The entries under $u_{1},u_{2},u_{3}$ are the
        primal surpluses $0$, $6$, $0$: the first and third requirements
        are met exactly ($1+3=4$, $2+3=5$), the second is over-fulfilled
        by $6$ ($1+9=10\ge 4$).

        (d)~Complementary slackness: the second requirement has surplus
        $6$, so $u_{2}=0$; $x_{1},x_{2}>0$ make both dual constraints
        equations: $u_{1}+2u_{3}=3$ and $u_{1}+u_{3}=2$, hence
        $u_{3}=1$, $u_{1}=1$. Both nonnegative, and
        $W=4+5=9=C^{*}$. $\checkmark$

  \item (a)~Sign rules for a max primal: labour and machine ($\le$) give
        $y_{1},y_{2}\ge 0$; the contract ($\ge$) gives $y_{3}\le 0$; the
        variables $x_{1},x_{2}\ge 0$ give $\ge$ dual constraints:
        \[
          \min\; W=10y_{1}+16y_{2}+6y_{3}
          \quad\text{s.t.}\quad
          y_{1}+2y_{2}\ge 5,\qquad y_{1}+y_{2}+y_{3}\ge 4,
        \]
        with $y_{1},y_{2}\ge 0$, $y_{3}\le 0$.

        (b)~At $x^{*}=(4,6)$: labour $4+6=10$ (binds), machine
        $8+6=14<16$ (slack $2$, so $y_{2}=0$), contract $x_{2}=6$ (binds).
        Both variables are positive, so both dual constraints are
        equations:
        \[
          y_{1}+2\cdot0=5\;\Rightarrow\;y_{1}=5,\qquad
          5+0+y_{3}=4\;\Rightarrow\;y_{3}=-1 .
        \]
        Signs: $y_{1}=5\ge 0$, $y_{2}=0$, $y_{3}=-1\le 0$, all as
        required, and both dual constraints hold. Values:
        $W=50+0-6=44$ and $Z^{*}=20+24=44$. $\checkmark$

        (c)~$y_{3}=-1$: every additional unit the contract demands lowers
        the maximal profit by $1$. Product~2 earns $4$ per unit but uses a
        labour hour that is worth $y_{1}=5$ in product~1, so forcing one
        more unit of it costs $5-4=1$. Check with a contract of $7$:
        labour and contract binding give $x=(3,7)$, machine $6+7=13\le 16$,
        and $Z=15+28=43=44-1$. $\checkmark$

        (d)~Written as $-x_{2}\le -6$, the row is of type $\le$ and gets a
        dual variable $\ge 0$; it is the old one with the sign flipped,
        $+1$. Read: raising the right-hand side $-6$ by one unit (to $-5$,
        a contract of $5$ units) raises $Z^{*}$ by $1$.

  \item (a)~Sign rules for a min primal: the $\ge$ rows give
        $u_{1},u_{3}\ge 0$, the $\le$ row gives $u_{2}\le 0$; the
        variables $x_{j}\ge 0$ give $\le$ dual constraints:
        \[
          \max\; W=8u_{1}+4u_{2}-10u_{3}
          \quad\text{s.t.}\quad
          u_{1}+u_{2}+u_{3}\le 2,\qquad 2u_{1}-3u_{3}\le 5,
        \]
        with $u_{1}\ge 0$, $u_{2}\le 0$, $u_{3}\ge 0$.

        (b)~At $x^{*}=(4,2)$: $4+4=8$ (binds), $x_{1}=4$ (binds),
        $4-6=-2>-10$ (surplus $8$, so $u_{3}=0$). Both variables are
        positive, so both dual constraints are equations:
        \[
          2u_{1}=5\;\Rightarrow\;u_{1}=\tfrac52,\qquad
          \tfrac52+u_{2}=2\;\Rightarrow\;u_{2}=-\tfrac12 .
        \]
        Signs: $u_{1}\ge 0$, $u_{2}\le 0$, $u_{3}=0$, all as required.
        Values: $W=8\cdot\tfrac52+4\cdot(-\tfrac12)-0=20-2=18$ and
        $C^{*}=8+10=18$. $\checkmark$

        (c)~$u_{1}=\tfrac52$: one more unit of the requirement raises the
        minimal cost by $2.5$ (it is met at the margin by $x_{2}$, which
        delivers $2$ units for $5$). $u_{2}=-\tfrac12$: one more unit of
        capacity for input~1 lowers the cost by $0.5$, because a unit of
        $x_{1}$ replaces half a unit of $x_{2}$: it saves $2.5$ and costs
        $2$. $u_{3}=0$: the third row has surplus, so its (negative)
        right-hand side is irrelevant at the margin. Check with
        $x_{1}\le 5$: $x=(5,1.5)$ satisfies $5+3=8$ and $5-4.5\ge -10$,
        and $C=10+7.5=17.5=18-0.5$. $\checkmark$

  \item (a)~$\mathcal{L}=3x_{1}+2x_{2}+\lambda_{1}(4-x_{1}-x_{2})
        +\lambda_{2}(6-2x_{1}-x_{2})$.

        (b)~With $x\ge 0$:
        \[
          \begin{aligned}
            &\pd{\mathcal{L}}{x_{1}}=3-\lambda_{1}-2\lambda_{2}\le 0,\quad
             x_{1}\ge 0,\quad x_{1}(3-\lambda_{1}-2\lambda_{2})=0,\\
            &\pd{\mathcal{L}}{x_{2}}=2-\lambda_{1}-\lambda_{2}\le 0,\quad
             x_{2}\ge 0,\quad x_{2}(2-\lambda_{1}-\lambda_{2})=0,\\
            &4-x_{1}-x_{2}\ge 0,\quad \lambda_{1}\ge 0,\quad
             \lambda_{1}(4-x_{1}-x_{2})=0,\\
            &6-2x_{1}-x_{2}\ge 0,\quad \lambda_{2}\ge 0,\quad
             \lambda_{2}(6-2x_{1}-x_{2})=0 .
          \end{aligned}
        \]
        (c)~At $x=(2,2)$, $\lambda=(1,1)$: $3-1-2=0$ and $2-1-1=0$ (so
        both products vanish); $4-4=0$ and $6-6=0$ (both rows bind, so
        both products vanish); all signs hold. Every condition is
        satisfied, and $Z=6+4=10$.

        (d)~$\partial\mathcal{L}/\partial x_{j}\le 0$ is
        $\lambda_{1}+2\lambda_{2}\ge 3$, $\lambda_{1}+\lambda_{2}\ge 2$,
        the dual constraints of Solution~1 with $\lambda$ in place of $y$;
        $\lambda\ge 0$ is the dual sign condition; the rows
        $\partial\mathcal{L}/\partial\lambda_{i}\ge 0$ are primal
        feasibility; the four product conditions are complementary
        slackness. So the multipliers are the dual solution
        $y^{*}=(1,1)$, i.e.\ the shadow prices. The objective is linear
        (concave) and the constraints are linear, so the Kuhn--Tucker
        conditions are sufficient: $(2,2)$ is optimal, as Solution~14
        found from the tableau.

  \item With $x_{j}\ge 0$ written as $-x_{j}\le 0$ and multiplier
        $\mu_{j}\ge 0$:
        \[
          \mathcal{L}=11x_{1}+14x_{2}+13x_{3}
          +\lambda_{1}(60-3x_{1}-2x_{2}-4x_{3})
          +\lambda_{2}(60-x_{1}-4x_{2}-x_{3})
          +\mu_{1}x_{1}+\mu_{2}x_{2}+\mu_{3}x_{3} .
        \]
        Stationarity at $\lambda=(3,2)$, $\mu=(0,0,1)$:
        \[
          \begin{aligned}
            \pd{\mathcal{L}}{x_{1}}&=11-3\lambda_{1}-\lambda_{2}+\mu_{1}
              =11-9-2+0=0,\\
            \pd{\mathcal{L}}{x_{2}}&=14-2\lambda_{1}-4\lambda_{2}+\mu_{2}
              =14-6-8+0=0,\\
            \pd{\mathcal{L}}{x_{3}}&=13-4\lambda_{1}-\lambda_{2}+\mu_{3}
              =13-12-2+1=0 .
          \end{aligned}
        \]
        Primal feasibility at $x=(12,12,0)$: $36+24=60\le 60$,
        $12+48=60\le 60$, $x\ge 0$. Signs: $\lambda,\mu\ge 0$.
        Complementary slackness: both rows bind, so
        $\lambda_{i}(60-\dots)=0$; $\mu_{1}x_{1}=0\cdot12=0$,
        $\mu_{2}x_{2}=0\cdot12=0$, $\mu_{3}x_{3}=1\cdot0=0$. All conditions
        hold, so the point is optimal (with $Z=300$). Interpretation:
        $\mu_{3}=(4\lambda_{1}+\lambda_{2})-13=1$ is the dual slack $t_{3}$
        of Additional Solution~9: the multiplier of the sign constraint
        $x_{3}\ge 0$ measures how much the profit would drop per unit if
        product~3 were forced into the plan, i.e.\ the margin it lacks.

  \item (a)~\textbf{False.} $t_{2}=4$: the margin must rise by more than
        $4$, to above $10=2y_{1}^{*}$. At $9$ the product still uses
        $2$ units of resource~A worth $10$ and earns only $9$.

        (b)~\textbf{True.} $t_{1}=1$, and a rise of $1.5>1$ lifts the
        margin to $5.5$, above the opportunity cost $1\cdot y_{1}^{*}=5$:
        the product now prices out positive ($5.5-5=0.5>0$), so the
        current plan is no longer optimal and product~1 enters.

        (c)~\textbf{False.} $t_{3}=0$ only says that product~3 breaks even
        at the current shadow prices, which it must, since it is made. Its
        margin sets the price $y_{1}^{*}=c_{3}$; if $c_{3}$ falls below
        $4$, product~1 earns more per unit of resource~A ($4$ against
        $c_{3}$), and the plan switches to product~1.

        (d)~\textbf{True.} The $z$-row entry under $x_{j}$ in the end
        tableau is $(A^{\top}y^{*})_{j}-c_{j}=t_{j}$, and the reduced cost
        is $c_{j}-(A^{\top}y^{*})_{j}=-t_{j}$; for basic variables both
        are $0$.
\end{solutions}
